% !TeX program = xelatex
\documentclass{book}
\usepackage[left=1in,right=1.5in,top=1in,bottom=1in]{geometry}
\usepackage{fontspec}
\IfFontExistsTF{Songti SC}
    {\setmainfont{Songti SC}}
    {\setmainfont{FandolSong-Regular}}
\usepackage{hyperref}

\hypersetup{
    colorlinks=true,
    linkcolor=black,
    urlcolor=cyan,
}

\title{簡明法語教程下冊}
\author{Junzhang Luo}

\begin{document}

\maketitle
\tableofcontents

\chapter{詞彙}

    \section{Double Meaning}
        \begin{enumerate}
            \item pauvre means ``可憐的'' before noun, after noun means ``貧窮的''
        \end{enumerate}

    \section{強調語氣}
        \begin{itemize}
            \item là 在句中起強調語氣的作用: C'est là votre erreur.
            \item à quel point = combien 表示`多麼': 
                \begin{center}
                    On voit bien à quel point il est attaché à sa patrie. 人們清楚地看到他是多麼熱愛自己的祖國
                \end{center}
            \item à [elle] seule 修飾 cette somme, 起強調作用，意思為僅此$\cdots$
                \begin{center}
                    Cette somme, à elle seule
                \end{center}
        \end{itemize}

    \section{連詞}
        \begin{itemize}
            \item et pourtant: 連結兩個意思對立的從句
            \item non seulement 不僅
            \item non seulement $\cdots$ mais $\cdots$ 不僅$\cdots$而且$\cdots$
            \item alors que: J'étais au travail, alors que vous dormiez encore.
            \item contraire à 與 $\cdots$ 相反, 對 $\cdots$ 不利
            \item outre + \textit{qch} 除$\cdots$之外，不僅:
                \begin{center}
                    Outre la patience, il manque de courage. 他不僅沒有耐心，還缺乏勇氣
                \end{center}
            \item soit 即，等於說: soit un butin de près de 1 million de dollars
            \item aboutir à \textit{qch} 通向，導致，以$\cdots$為結果
            \item si \dots que \dots 如此$\cdots$以至於$\cdots$
            \item n'en \dots pas moins (nevertheless)
            \item devoir 可以表示``也許，大概''; Elle ne devait pas avoir de parapluie 她大概沒有雨傘
            \item trop \dots pour que \dots 太$\cdots$以至於$\cdots$
            \item à \dots près 除了$\cdots$
        \end{itemize}

    \section{``ne'' Related}
        \begin{itemize}
            \item ne \dots plus que 只剩下，只有
            \item ne faire que + \textit{inf}
            \item ne \dots ni \dots ni 既不 $\cdots$ 也不 $\cdots$
        \end{itemize}

    \section{Number Related}
        \begin{itemize}
            \item des dizaines de milliers 數萬
            \item des centaines de milliers 數十萬
            \item mesurer $\cdots$ de long (large) 長(寬)度為
            \item y compris, non compris 包括，不包括 (若放在所修飾名詞後面, agrees to gender and number)
            \item à peine 剛剛，才，勉強
            \item autant que 於$\cdots$相等
            \item presqu'autant que 與$\cdots$相差無幾
            \item tendre à 有$\cdots$的趨勢，以$\cdots$為目的
        \end{itemize}

        \subsection{確切數量}
            \begin{enumerate}
                \item Combien d'étudiants y a-t-il dans la classe? Il y en a 20 (étudiants).
                \item Quelle est \textbf{la longueur} de cette pièce? Elle a \textbf{cinq mètres} de long.
                \item Quelle est \textbf{la largeur} de cette pièce? Elle est large de 1.5 \textbf{centimètres}.
                \item De quelle \textbf{hauteur} est la tour Eiffel? Elle est haut de 320 \textbf{mètres}.
                \item Quel est \textbf{le poids} de cette valise? Elle pèse \textbf{10 kilos}.
            \end{enumerate}
        \subsection{非確切數量}
            \begin{enumerate}
                \item 使用部分冠詞
                    \begin{center}
                        J'ai \textbf{de l'}argent. Il veut \textbf{du} café. La grand-mère a acheté \textbf{des} légumes au marché.
                    \end{center}
                    否定用de: Nous n'avons pas \textbf{de} fruits.
                \item 數量副詞 + de: \textbf{assez de} temps, \textbf{un peu de} pain, \textbf{beaucoup de} jeunes
            \end{enumerate}
        \subsection{小數量}
            \begin{enumerate}
                \item \textbf{Certains} élèves ont mal compris la question.
                \item Nous sommes \textbf{plusieurs} à dîner.
                \item \textbf{Plus d'un} restaurant est fermé à 20 heures.
                \item J'ai \textbf{quelques} remarques à faire.
                \item Il en manque \textbf{quelques-uns}.
                \item Je n'ai \textbf{guère} le temps de faire la cuisine.
            \end{enumerate}
        \subsection{大數量}
            \begin{enumerate}
                \item \textbf{Un grand nombre de} jeunes couples ne veulent pas avoir des enfants.
                \item \textbf{Nombre d'}étudiants ont assisté à la réunion.
                \item Il invoque \textbf{une foule de} prétextes pour ne pas venir. (A bunch of)
                \item Elle trouve toujours \textbf{un tas d'}excuses pour e pas venir au rendez-vous. (a lot of)
                \item \textbf{Beaucoup d'}enfants de 2 à 6 ans vont à l'école maternelle.
                \item Ces immigrés ont rencontré \textbf{pas mal de} problèmes depuis leur arrivée en Europe. (quite a few)
                \item En cette saison, il y a \textbf{trop de} monde sur la plage.
            \end{enumerate}
        \subsection{約略數}
            \begin{enumerate}
                \item Il nous faut attendre \textbf{environ} 20 minutes.
                \item Elle gagne \textbf{autour de} 3000 euros par mois.
                \item Il y a \textbf{près d'}une demi-heure que je suis ici.
                \item Nous nous retrouverons \textbf{vers} 6 heures.
                \item Les professeurs ont \textbf{à peu près} 8 heures de cours par semaine.
                \item C'est une femme \textbf{entre} 30 et 35 ans.
                \item Il a \textbf{à peine} 20 ans. (barely)
                \item Les candidats sont \textbf{au nombre de} 200 environ.
            \end{enumerate}
            

    \section{Expression Du Temps}
        \begin{enumerate}
            \item avant 在 $\cdots$ 之前
            \item après 在 $\cdots$ 之後
            \item il y a \dots $\cdots$以前
            \item depuis \dots $\cdots$以來
            \item ça fait \dots que 已經$\cdots$
            \item dans \dots $\cdots$ 之後
            \item en un temps donné 在一定的時間內
            \item le plus souvent 往往，通常
            \item travailler à plein temps (mi-temps) 全日制工作, 半日制
            \item beau 表示某一天，某一個: un beau matin 一天早晨
        \end{enumerate}

        \subsection{時間 Le Temps}
            \begin{enumerate}
                \item 確切時間
                    \begin{itemize}
                        \item à trois heures
                        \item le 10 janvier
                        \item lundi prochain
                        \item au mois de février (en février)
                        \item au printemps, en éte, en automne, en hiver
                        \item il y a 
                        \item dans
                        \item à [time] à partir du
                    \end{itemize}
                \item 非確指時間
                    \begin{itemize}
                        \item début août (fin juin) 
                        \item à la fin du mois
                        \item dans les années 50
                    \end{itemize}
                \item 重複
                    \begin{itemize}
                        \item le [weekday]
                        \item une fois par an
                        \item toutes les 3 minutes
                    \end{itemize}
                \item 時間副詞
                    \begin{itemize}
                        \item aujourd'hui
                        \item tout à l'heure (just now)
                        \item parfois
                        \item tout de suite (right away)
                        \item de temps en temps
                        \item tout le temps 
                        \item aussitôt (right away)
                        \item aussitôt après (immediately after)
                    \end{itemize}
            \end{enumerate}
        
        \subsection{期限 La Durée}
            \begin{enumerate}
                \item 與 ``說話時'' 有關
                    \begin{itemize}
                        \item J'attends le car \textbf{depuis} une demi-heure.
                        \item \textbf{Il y a} trois jours \textbf{que} je suis ici.
                        \item \textbf{Voilà} bientôt deux ans \textbf{qu'}ils va à l'école. (it's been almost)
                        \item \textbf{Ça fait} une semaine \textbf{qu'}ils ne se sont pas revus. (it's been )
                    \end{itemize}
                \item 與 ``說話時'' 無關
                    \begin{itemize}
                        \item Il a été absent \textbf{pendant} deux semaines.
                        \item Il est parti aux États-Unis \textbf{pour} deux ans.
                        \item Ces armes ont été inventées \textbf{durant} la première guerre mondiale.
                        \item Ce bâtiment a été construit \textbf{en} 10 jours.
                        \item \textbf{En l'espace de} 2 heures vous serez à Guangzhou.
                    \end{itemize}
            \end{enumerate}

        \subsection{時間關係 Les Rapports de Temps}
            \begin{enumerate}
                \item La simulatanéité 同時性
                    \begin{itemize}
                        \item Quand
                        \item Lorsque (while)
                        \item Pendant que (while)
                        \item alors que (while)
                        \item à mesure que (as)
                        \item chaque fois (each time)
                        \item tant que (while)
                        \item tabdis que (while)
                    \end{itemize}
                \item L'antériorité 先時性
                    \begin{itemize}
                        \item avant que + subjonctif
                        \item avant de + inf
                        \item en attendant que + subjonctif
                        \item jusqu'à ce qu + subjonctif
                        \item jusqu'au moment où + indicatif
                    \end{itemize}
                \item La postériorité 后时性
                    \begin{itemize}
                        \item après que + indicatif
                        \item après + infinitif passé
                        \item quand + forme compossée: \textbf{Quand j'ai fini} mes études, je me suis engagé dans l'armée.
                        \item dès que (une fois, que, aussitôt que) + form composée (as soon as, once, that, as soon as)
                        \item depuis que + indicatif (since)
                    \end{itemize}
            \end{enumerate}
    
    \section{Position}
        \begin{enumerate}
            \item sur
            \item sous
            \item devant
            \item derrière
            \item dans
            \item entre
            \item en face de
            \item au milieu de 在$\cdots$中間
            \item au centre de 在$\cdots$中央
            \item en dehors de
            \item à côté de 旁邊
            \item près de 附近
            \item ``mi-'' 前綴詞，表示一半，正中: mi-nez 鼻子中部
        \end{enumerate}
    
    \section{Comparison}
        La supériorité
        \begin{itemize}
            \item être inférieur (supérieur) à 低於，高於，優於
            \item moindre 与定冠词(le) 连用, 構成 petit 最高級
            \item valoir mieux (que) 比$\cdots$好，勝於
            \item plus \dots que, plus de \dots que
            \item de plus en plus
            \item davantage
            \item plus \dots plus
            \item d'autant plus \dots que
        \end{itemize}
        L'inféroirité
        \begin{itemize}
            \item moins \dots que, moins de \dots que
            \item de moins en moins
            \item moins \dots moins
            \item d'autant moins \dots que
        \end{itemize}
        \subsection{L'égalité et la ressemblance 相等與相似}
            \begin{enumerate}
                \item aussi + adj/adv + que
                \item autant + que
                \item autant de + nom + que
                \item comme, comme si
                \item le même, selon, suivant, conformément à
                \item tel, tel que
                \item ainsi que, de même que
                \item autant \dots autant \dots
                \item on dirait que
            \end{enumerate}
    
    \section{泛指 ``Il'' 相關}
        \begin{itemize}
            \item il est question de 問題在於，事關
            \item il vaut mieux + \textit{inf} 最好是$\cdots$
            \item vaincu (comme il était vaincu) 表示原因
            \item il me paraît impossible de 依我看
            \item il est aisé de 容易$\cdots$，不難$\cdots$
            \item il est permis à \textit{qn} de faire \textit{qch} 允許某人做某事
            \item il arrive que + subjonctif 有時發生$\cdots$, 有時會$\cdots$
            \item il n'est pas admissible que + subjonctif 無法接受
        \end{itemize}
    
    \section{Mots et Expressions}

        \begin{itemize}
            \item avoir mal à 身体部位不适
            \item l'échapper belle 幸免
            \item à la sortie de \textit{loc. prép}: 在$\cdots$ 結束時，離開時
            \item en être 表示一件事情的進展程度: où en est votre travail? 您的工作進行的怎麼樣了?
            \item pris à 從某處獲得 (from prendre): les canons qu'on avait pris à l'ennemi.
            \item se rendre compte de \textit{qch} (que) 體會到，意識到: Je me rends compte que ces étudiants connaissent beaucoup de choses.
            \item quant à \textit{loc. prép} 至於，關於: Pars si tu veux; quant à moi, je reste.
            \item aux yeux de \textit{qn} 在$\cdots$眼裡， 依$\cdots$看來
            \item avoir l'air 好像，似乎 (adj. follows agree to gender and number)
            \item 為$\cdots$高興，自豪: (être) content/fier de
            \item accorder de l'importance à 對$\cdots$給予重視
            \item (être) orné de 裝飾著
            \item de choix \textit{loc. adj} 精選的，上乘的
                \begin{center}
                    L'illumination des monuments et des fontaines offre au regard un spectacle de choix.
                    \\
                    au regard (to the eye)
                \end{center}
            \item en avoir pour 為此需要$\cdots$
            \item se réjouir de \textit{qch} 喜悅，高興
            \item à la légère 草率的
            \item tirer un avantage de \textit{qch} 從$\cdots$得到好處
            \item être coupé de 與$\cdots$隔絕
            \item suivre l'exemple de 向$\cdots$學習
            \item se terminer par 以$\cdots$為結束
            \item en fonction de 根據$\cdots$, 依照$\cdots$
            \item le plus souvent 往往
            \item être chargé de 負責$\cdots$
            \item insister sur \textit{qch} 強調某事
            \item se méfier 懷疑，當心
            \item s'adapter à 適應$\cdots$
            \item (être) attaché à \textit{qch} 對$\cdots$喜愛，依戀
            \item estimer que 認為
            \item être suivi de 在$\cdots$之後: Son discours a été suivi d'un long applaudissement.
            \item sous les effets de 受$\cdots$的作用
            \item se passer de \textit{qch} 免去，省掉
            \item (être) doté de 被賦予$\cdots$，具有$\cdots$
            \item (être) indispensable à 對$\cdots$必不可少的
            \item se libérer de 擺脫$\cdots$
            \item tirer \textit{qn} par 拉住某人身體的一部位
            \item s'introduire 進入，傳入
            \item avoir recours à \textit{qch} 求助於，依靠$\cdots$
            \item accorder de l'attention à \textit{qch} 關心某事
            \item $\star$ à la main \textit{qch} 作方式狀語, 不會說 sa main verb. \textit{qch}
            \item couper la parole à \textit{qn} 打斷某人的話，使某人說不出話
            \item être doué(e) de 具有，富有$\cdots$
            \item avoir le sentiment que 覺得，感到$\cdots$
            \item être méchant avec \textit{qn} 對某人凶惡
            \item avec sursis 緩期執行
            \item transformer \textit{qch} en \dots 把$\cdots$轉變成
            \item appeler \textit{qch} 稱$\cdots$為$\cdots$
            \item à nouveau 再一次，重新
            \item n'importe (qui, quoi, comment, où, quand, lequel etc.) 不論$\cdots$
            \item à compter de 从$\cdots$算起
            \item de tout repos 可靠的，令人放心的
            \item être au courant 了解情況的，熟知$\cdots$的
            \item pardonner 一般用否定式，無可挽回的，致命的
                \begin{center}
                    C'est une erreur qui ne pardonne pas. 這是一個無法挽回的錯誤
                \end{center}
            \item en activité 在職的，起作用的
            \item se laisser + \textit{inf} 被$\cdots$，任憑$\cdots$
            \item avant l'âge 未到年紀
            \item savoir tout sur 精通$\cdots$
            \item ne savoir rien de 對$\cdots$一無所知
            \item mettre en œuvre 使用，發揮
        \end{itemize}

        \subsection{faire}
            faire + 帶部分冠詞(de)的名詞: 做$\cdots$, 從事$\cdots$
            \begin{itemize}
                \item faire de \textit{qch} 使用，安排
                \item faire signe à \textit{qn} 像某人示意
                \item faire partie de 屬於$\cdots$
                \item se faire prendre 被抓住，被查出
                \item faire loi 具有法律效力，起決定性的作用
            \end{itemize}

        \subsection{arriver}
            \begin{itemize}
                \item \textit{qch}. arrive à \textit{qn} 發生，來臨
                \item v. \textit{impers}: Il lui est arrivé un malheur 他遭遇不幸
            \end{itemize}

        \subsection{sembler}
            \begin{itemize}
                \item sembler à \textit{qn} 在某人看來，使某人覺得
                \item v. \textit{impers}: 在某人看來，似乎 
                    \begin{itemize}
                        \item Il semble qu'il il fait plus chaud aujourd'hui qu'hier.
                        \item Il me semble utile de faire ces exercices.
                    \end{itemize}
            \end{itemize}

        \subsection{essayer}
            \begin{itemize}
                \item essayer de 盡力，力圖
            \end{itemize}

        \subsection{travailler}
            \begin{itemize}
                \item travailler à 積極從事，竭力
            \end{itemize}

        \subsection{expliquer}
            \begin{itemize}
                \item expliquer \textit{qch} à \textit{qn} 向某人解釋某事
            \end{itemize}

        \subsection{proposer}
            \begin{itemize}
                \item propser (à \textit{qn}) de faire \textit{qch} 建議某人做某事: Je vous propose de vous conduire à la gare.
            \end{itemize}
        
        \subsection{remettre}
            \begin{itemize}
                \item remettre \textit{qch} à \textit{qn} 交給某人某物
            \end{itemize}

        \subsection{penser}
            que penser de \textit{qch} 對$\cdots$怎麼看: Que pensez-vous de ce film? 您覺得這部電影怎麼樣?
        
        \subsection{rendre}
            \begin{itemize}
                \item se rendre 投降
            \end{itemize}

        \subsection{apercevoir}
            看見，注意到
            \begin{itemize}
                \item s'apercevoir de \textit{qch} (que)
                    \begin{center}
                        Ne s'aperçoit-il jamais que ses auditeurs sont lassés de ses discours? \\
                        他就沒有發覺聽眾對他的演講不感興趣嗎?
                    \end{center}
            \end{itemize}

        \subsection{apprendre}
            \begin{itemize}
                \item apprendre à + \textit{inf} 學習做$\cdots$
                \item apprendre à \textit{qn} à + \textit{inf} 教某人，讓$\cdots$知道
            \end{itemize}

        \subsection{renseigner}
            \begin{itemize}
                \item renseigner \textit{qn} 告訴
                \item se renseigner 打聽情況
            \end{itemize}   
        
        \subsection{relever}
            扶起，升起
            \begin{itemize}
                \item se relever 重新站起，升高: À partir de ce point, le terrain se relève brusquement (突然).
            \end{itemize}
        
        \subsection{peser}
            稱，考慮，重(weight)
            \begin{itemize}
                \item peser sur 使有重覆之感，影響: La mort subite de son père va peser sur sa décision.
            \end{itemize}
        
        \subsection{regarder}
            \begin{itemize}
                \item regarder à \textit{qch} 考慮，注意
            \end{itemize}

        \subsection{servir}
            \begin{itemize}
                \item se servir de 使用，利用
                \item servir à 用於
            \end{itemize}

        \subsection{tarder}
            \begin{itemize}
                \item tarder à + \textit{inf} 延遲，遲緩
            \end{itemize}

        \subsection{voir}
            \begin{itemize}
                \item voir bien que 看出，發現
                    \begin{itemize}
                        \item Je vois bien que vous êtes malade. 
                    \end{itemize}
            \end{itemize}

        \subsection{revenir}
            再來，回來
            \begin{itemize}
                \item revenir à \textit{qch} 回到
                \item revenir de \textit{qch} 恢復，擺脫
            \end{itemize}

        \subsection{considérer}
            \begin{itemize}
                \item considérer que 認為
                \item considérer \textit{qch(qn)} comme 視作
            \end{itemize}

        \subsection{tomber}
            \begin{itemize}
                \item tomber malade 病倒
                \item tomber amoureux (de \textit{qn}) 鍾情與某人
                \item tomber d'accord 達成一致意見
            \end{itemize}
        
        \subsection{souffrir}
            \begin{itemize}
                \item souffrir de 受$\cdots$痛苦
            \end{itemize}

        \subsection{parvenir}
            achieve
            \begin{itemize}
                \item parvenir à + \textit{inf} 能夠，終於
            \end{itemize}

        \subsection{manquer}
            \begin{itemize}
                \item manquer \textit{qch} 錯過
                \item manquer de \textit{qch} 缺乏
                \item manquer à \textit{qn} 缺少，不足，失去
            \end{itemize}

        \subsection{obliger}
            \begin{itemize}
                \item obliger \textit{qn} 使承擔義務
                \item obliger \textit{qn} à + \textit{inf} 強迫，迫使
                \item être obligé de + \textit{inf} 必須，被迫
            \end{itemize}

    \section{相似詞}
        \begin{itemize}
            \item reconnaître 认出，辨认出, connaître 認識，了解, savoir 知道一個事實，會做某事
        \end{itemize}

\chapter{語法}
    \section{中性代詞 Le Pronom Neutre ``le''}
        無性數變化
        \begin{enumerate}
            \item 作直接賓語，代替一個不定時動詞: 
                \begin{center}
                    Partez, puisque vous le voulez = puisque vous voulez partir \\
                    既然您要走，您就走吧
                \end{center}
            \item 昨表語，代替身份職業名詞:
                \begin{center}
                    Est-ce que vous êtes journaliste? Oui, je \textbf{le} suis.
                \end{center}
            \item 代替形容詞: Êtes-vous fatigués? Oui, nous \textbf{le} sommes.
            \item 代替句子: 
                \begin{center}
                    Vous êtes satisifait, je \textbf{le} vois.
                \end{center}
        \end{enumerate}

    \section{主有代詞 Le Pronom Possessif}
        \begin{center}
            \begin{tabular}{|c|c|c|c|}
                \hline
                le mien & le mienne & le miens & le miennes \\
                \hline
                le tien & le tienne & le tiens & le tiennes \\
                \hline
                le sien & le sienne & le siens & le siennes \\
                \hline
                le nôtre & la nôtre & les nôtres & les nôtres \\
                \hline
                le vôtre & la vôtre & les vôtres & les vôtres \\
                \hline
                le leur & la leur & les leurs & les leurs \\
                \hline
            \end{tabular}
        \end{center}

        \subsection{用法}
            \begin{enumerate}
                \item 代替主有形容詞(mon, ton, son etc.) 加名詞，避免名詞重複
                    \begin{center}
                        Ma chambre est plus grande que la vôtre = votre chambre
                    \end{center}
                \item 陽性複數形式可以不代替任何名詞，表示：家人，親友，團體成員等
                    \begin{center}
                        Ils sont très heureux de retrouver les leurs. 他們與親友團聚，非常高興
                        \\
                        Ne dites pas cela à eux, ils ne sont pas des nôtres. 別把這事告訴他們，他們和我們不是一夥的
                    \end{center}
            \end{enumerate}
        
        \subsection{Remark}
            \begin{enumerate}
                \item 若定冠詞前有介詞 à or de, 變成 au mien, du tien, des leurs etc.
                \item 與被佔有者的性，數變化 (目標本身)
            \end{enumerate}

    \section{复合关系代词 Le Pronom Relatif Composé}
        lequel 的变化，性数一致，à,de 变化
        \begin{tabular}{|c|c|c|c|}
            \hline
            lequel & laquelle & lesquels & lesquelles \\
            \hline
            auquel & à laquelle & auxquels & auxquelles \\
            \hline
            duquel & de laquelle & desquels & desquelles \\
            \hline
        \end{tabular}

        \subsection{用法}
            \begin{enumerate}
                \item 關係從句中作主語，代替關係代詞qui (more in written language)
                    \begin{center}
                        J'ai rencontré la sœur de Pierre, laquelle va partir pour le Canada. (laquelle = la sœur de Pierre)
                    \end{center}
                \item à, de 或介詞短語連用, 在關係從句中作間接賓語或狀語(若先行詞之人，也可以用 à qui 形式)
                    \begin{center}
                        Il m'a posé des questions auxquelles je n'ai pas pu répondre\\
                        (auxquelles = à ces questions, 作 répondre 的間接賓語: Je n'ai pas pu répondre à ces questions.)
                        \\
                        Voilà monsieur Dupont auquel (ou: à qui) j'écris souvent. \\
                        (auquel = à monsieur Dupoint: J'écris souvent à monsieur Dupont.)
                    \end{center}
                \item 和其他介詞或介詞短語連用, 在關係從句中作狀語
                    \begin{center}
                        Les étudiants ont tenu une réunion au cours de laquelle ils ont élu le président de leur association. 
                        \\
                        (au cours de laquelle = au cours de la réunion)
                    \end{center}
            \end{enumerate}

        \subsection{Remark}
            \begin{enumerate}
                \item Do not mix up with 复合疑问词: lequel, laquelle etc.
            \end{enumerate}

    \section{不定式 L'infinitif}
        \subsection{用法}
            \begin{enumerate}
                \item 作主語或邏輯主語：
                    \begin{center}
                        Partir est triste = Il est triste de partir. 離去是令人傷心的
                        \\
                        Cela me plaît d'aller au théâtre = Aller au théâtre me plaît 我喜歡看戲
                    \end{center}
                \item 作賓語: 取決於前面的動詞，有時需要介詞 à, de 引導
                    \begin{enumerate}
                        \item 直接跟動詞後: J'aime faire du sport 
                        \item 由 à 引導: Nous commençons \textit{à} comprendre.
                        \item 由 de 引導: Elle lui demande \textit{de} venir.
                    \end{enumerate}
                \item 在一些介詞或短語後使用: être en train de etc.
                \item 構成省略疑問句: Comment faire? Pourquoi partir si tôt?
            \end{enumerate}

        \subsection{時態}
            \begin{enumerate}
                \item 不定式現在時 (l'infinitif présent): 表示與主要動詞的動作同時或在其之後發生的動作
                    \begin{center}
                        Je veux \textit{lire} ce roman. Je lui ai demandé de \textit{venir} me voir.
                    \end{center}
                \item 不定式過去時 (l'infinitif passé): être, avoir + 過去分詞, 表示與主要動詞比較已經完成的動作
                    \begin{center}
                        Je me souviens d'\textit{avoir lu} ce roman.
                        \\
                        Il regrette de vous \textit{avoir fait} attendre si longtemps.
                    \end{center}
            \end{enumerate}
    
    \section{過去分詞 Le Participe Passé}
        Agrees in gender, number when
        \begin{enumerate}
            \item avoir, 且直接賓語在動詞前面時
            \item 被動意義，純粹的代動詞 (去掉se etc. 後動詞意思不一樣)
            \item 過去分詞獨立使用，作表語 (subject complement, 狀態)，定語 (adjective modifier, like ``a letter written in'') 或同位語時
                \begin{center}
                    Nous accueillons des amis venus (= qui sont venus) de différents pays du monde.
                \end{center}
        \end{enumerate}

    \section{復合過去分詞 Le Participe Passé Composé}
        être, avoir 的現在分詞 (ayamt, étant) + 動詞過去分詞, 表示在主要動詞之前就已經完成的動作
            \begin{itemize}
                \item 主動態: demander $\to$ ayant demandé
                \item 被動態: faire $\to$ ayant été fait
            \end{itemize}
            \begin{center}
                Bernard, ayant pris la parole, a quitté la salle (ayant pris = qui avait pris) 
                \\
                Ayant passé les examens, les étudiants sont rentrés chez eux (ayant passé = après avoir passé)
                \\
                (Étant) arrivés trop tard à la gare, ils ont manqué leur train (arrivés = comme ils étaient arrivés)
            \end{center}

    \section{à 用法}
        \begin{enumerate}
            \item 去向，地點: Nous allons à Beijing.
            \item 距離: La gare est à 100 mètres d'ici.
            \item 時間: Ce monument a été construit au $18^{e}$ siècle
            \item 方式: Il va à vélo.
            \item 所屬: Ce magnétophone est à moi.
            \item 引導名詞補語: C'est une machine à écrire.
            \item 引導形容詞補語: Il est fidèle à sa petite amie.
            \item 引導間接賓語: Je dis au revoir à Pierre.
        \end{enumerate}

    \section{de 用法}
        \begin{enumerate}
            \item 來源，從$\cdots$，自$\cdots$: Je vien de Paris.
            \item 分量，內容: un kilo de sucre. Je voudrais une tasse de café.
            \item 程度: Ma montre retarde de deux minutes. La production a augmenté de 5\%
            \item 方式: Il m'a fait signe de la tête.
            \item 引導名詞補語, $\cdots$的: C'est la moto de Pierre.
            \item 引導形容詞補語: Ce mur est haut de 3 mètres.
            \item 引導間接賓語: Il se souvient de son voyage en Chine.
            \item 引導不定式動詞: Je suis très content de vous revoir.
        \end{enumerate}

    \section{à de 連用}
        \begin{enumerate}
            \item de \dots à \dots 從$\cdots$到$\cdots$ (表示空間或時間)
                \begin{center}
                    le voyage de la Terre à la Lune
                \end{center}
            \item 分別引導直接賓語和間接賓語
                \begin{center}
                    Je demande à Paul de m'expliquer ce problème.
                \end{center}
        \end{enumerate}
    
    \section{en 用法}
        \begin{enumerate}
            \item 地點: Les étudiants sont en classe.
            \item 時間: Ils sont retournés à Naples en août.
            \item 狀態: La ville est en fête. 
            \item 在$\cdots$方面: Il est très fort en mathématiques.
            \item 方式: Je lui ai donné ce livre en cadeau.
            \item 材料: Nous habitons dans une maison en briques.
            \item 顏色: Ces photos sont en couleurs (en noir et blanc).
        \end{enumerate}
    \section{aussi ainsi}
        \begin{itemize}
            \item 以aussi, ainsi 等詞引導的陳述句，主語和謂語要倒裝
                \begin{center}
                    Aussi les végétaux doivent-ils absorber une grande quantité de liquide.
                \end{center}
        \end{itemize}
        
    \section{tout}
        \begin{enumerate}
            \item 作泛指形容詞：所有的，整個的 (agrees to gender and num)
            \item 作泛指代詞：全部，整個，大家 (agrees to gender and num)
            \item 作副詞：完全地，很，非常 (does not agree, unless the adj. follow 為陰性 and starts with consonant or silent `h')
        \end{enumerate}
    
    \section{avec}
        \begin{enumerate}
            \item 和$\cdots$一起
            \item 同，與
            \item 隨著，與$\cdots$同時
            \item 表示``方式'', 也可以用de: d'un ton calme
            \item 帶有，具有
        \end{enumerate}

    \section{dont}
        代替有介词de的(先行)名詞
        \begin{enumerate}
            \item 動詞的間接賓語: Voilà le docteur dont je vous ai parlé (dont = de ce docteur).
            \item 名詞補語: C'est un homme dont la vie est consacrée à la science (dont = de cet homme, la vie de cet homme).
            \item 形容詞補語: Il nous a parlé de son travail dont il est très satisfait (satisfait de son travail).
        \end{enumerate}
    
    \section{pour}
        \begin{enumerate}
            \item 目的地
            \item 時間
            \item 目的
            \item 對象，用途
            \item 依$\cdots$看，對$\cdots$來說
            \item 價格
            \item 作為，當作
        \end{enumerate}

    \section{comme}
        \begin{enumerate}
            \item 作爲，當作: 
                \begin{center}
                    Comme citoyen, j'ai le droit de voter. 作爲公民，我有選舉的權利
                    \\
                    Il travaille comme journaliste. 他做記者工作
                \end{center}
            \item 像，如同
            \item 當$\cdots$時
                \begin{center}
                    Le téléphone a sonné comme je sortais. 我正要出門電話鈴響了
                \end{center}
            \item 例如
            \item 由於，既然
            \item 多麼，怎樣(感嘆)
        \end{enumerate}

        \subsection{詞組}
            \begin{itemize}
                \item se considérer comme 自認為是
            \end{itemize}

    \section{par}
        \begin{enumerate}
            \item 表示地點: 經過，從
            \item ``每''
            \item 方式或方法: 用，以，通過
            \item 引導施動者補語: Ce raps a été préparé par mon père.
            \item 用在 commencer 和 finir 後，表示以$\cdots$開始，結束: Ils ont fini par trouver un appartement.
        \end{enumerate}

\chapter{時態}
    \section{配合}
        \begin{enumerate}
            \item 主語動詞現在時
            \item 主語動詞將來時: Je travaillerai
                \begin{itemize}
                    \item quand j'aurai 18 ans 從句與主句同時
                    \item quand j'aurai obtenu mon diplôme 從句先與主句
                \end{itemize}
            \item 主語動詞過去時: Je croyais ou J'ai cru
                \begin{itemize}
                    \item qu'il était malade 同時
                    \item qu'il partirait le lendemain 從句後於主句
                    \item qu'il était déjà parti 先與
                \end{itemize}
        \end{enumerate}
    
    \section{復合過去時 Passé Composé}

    \section{未完成過去時 Imparfait}

        \subsection{用法}
            過去重複，持續，背景性的狀態/動作

        \subsection{Conjugation}
            nous present stem 去掉 ``ons'', 加
            \begin{itemize}
                \item ais, ais, ait, ions, iez, aient
            \end{itemize}

    \section{簡單過去時 Passé Simple}
        
        \subsection{Conjugation}
            \begin{itemize}
                \item ai, as, a, âmes, âtes, èrent
                \item is, is, it, îmes, îtes, irent
                \item us, us, ut, ûmes, ûtes, urent
            \end{itemize}

        \subsection{Irregular}
            être (fut), avoir (eut), faire (fit), venir (vint), voir (vit), pouvoir (put), devoir (dut), vouloir (voulut), savoir (sut), prendre (prit), mettre (mit), dire (dit)
        
        \subsection{用法}
            某一確定時間內已完成的動作，書面語言，通常只使用第三人稱單數和複數, passé composé 的書面語
        
        \subsection{Remark}
            如果緊密相連的兩個動作，則要用先過去時(le passé antérieur, avoir/être 簡單過去時加動詞過去分詞), 通常用在 dès que, aussitôt que, à peine que (一$\cdots$ 就)等短語引導的時間狀語從句中:
            \begin{center}
                Dès qu'il \textbf{eut fini} son travail, il aida les autres. 
            \end{center}

    \section{愈過去時 Le plus-que-parfait}

        \subsection{構成}
            avoir être 未完成過去時(imparfait) + 過去分詞

        \subsection{用法}
            \begin{enumerate}
                \item 在過去某時前已發生，一般配合使用
                    \begin{center}
                        Il m'a demandé: << Quand Marie est-elle partie? >> \\
                        Il m'a demandé quand Marie était partie.
                    \end{center}
                \item 直接引語變成間接引語，除了根據句子改變人稱，時態，指示代詞 (ce, cet, cette, ces, ceci, cela, ça, ce)，指示形容詞(celui, celle, ceux, celles) 外，有時候還需改變時間副詞:
                    \begin{itemize}
                        \item maintenant, en ce moment $\to$ alors, à ce moment-là 現在此時，當時那時候
                        \item ce soir $\to$ ce soir-là
                        \item aujourd'hui $\to$ le jour même, ce jour-là
                        \item hier $\to$ la veille
                        \item il y a deux jours $\to$ deux jours plus tôt, dans deux jours $\to$ deux jorus plus tard
                    \end{itemize}
                    \begin{center}
                        Il dit: << Je suis très occupé maintenant. >> 
                        \\
                        Il a dit qu'il était très occupé à ce moment-là.
                    \end{center}
                    
            \end{enumerate}

    \section{現在分詞 Le Participe Présent}
        \subsection{構成}
            ons steam with ant
        \subsection{Irregular}
            ayant, étant, sachant
        \subsection{用法}
            \begin{enumerate}
                \item 關係代詞 qui+變位動詞
                    \begin{center}
                        On recrute des interprètes connaissant à la fois le français et l'anglais. \\
                        (connaissant = qui connaissent)
                    \end{center}
                \item 表示時間或原因的狀語從句:
                    \begin{center}
                        Ne sachant comment faire, elle téléphona à la police. 
                        \\
                        (Comme elle ne savit pas comment faire)
                    \end{center}
                \item 當形容詞使用 (agrees to gender and num)
                    \begin{center}
                        La situation est encourageante.
                    \end{center}
            \end{enumerate}

    \section{副動詞 Le Gérondif}
        en + 現在分詞
        \begin{enumerate}
            \item 作時間狀語，表示同時性，可以在 en 前面加 tout 來強調同時性
                \begin{center}
                    En allant à la poste, j'ai vu Pierre. Elles bavardaient tout en tricotant.
                \end{center}
            \item 方式狀語: C'est en nageant qu'on apprend à nager.
            \item 條件狀語: En faisant beaucoup \dots (= Si vous faites \dots)
        \end{enumerate}

    \section{條件式現在時 Le Conditionnel Présent}
        構成與過去將來時一樣簡單將來時詞根 + ais, ais, ait, ions, iez, aient
        
        \subsection{用法}
            \begin{enumerate}
                \item 婉轉的請求，建議，推測
                \item 主從復合句，表示假設或可能，配合: 從句(si + 直述式未完成過去時)，主句(條件現在時)
                \begin{enumerate}
                    \item 表示與現在事實相反的假設
                        \begin{center}
                            Je vous aiderais, si je le pouvais. S'il avait le temps, il apprendrait aussi l'italien.
                        \end{center}
                    \item 表示將來可能(較小)發生的動作
                        \begin{center}
                            Si elle venait ici demain, elle pourrait répondre à ces questions.
                        \end{center}
                    \item 若可能性較大，用陳述式表示條件和結果 (主: 簡單將來時，從: si + 直陳式現在時)
                        \begin{center}
                            S'il pleut demain, nous ne sortirons pas.
                        \end{center}
                \end{enumerate}
            \end{enumerate}

    \section{簡單將來時}
        ai, as, a, ons, ez, ont

    \section{虛擬現在時 Le Subjonctif Présent}
        ent stem with e, es, e, ions, iez, ent
        \subsection{Irregular}
            \begin{itemize}
                \item avoir: aie, aies, ait, ayons, ayez, aient
                \item être: sois, sois, soit, soyons, soyez, soient
                \item aller: aille, allions
                \item faire: fasse
                \item pouvoir: puiss-
                \item savoir: sach-
                \item vouloir: veuille, voulions
                \item falloir: faille
                \item pleuvoir: pleuve
            \end{itemize}
        \subsection{用法}
            \begin{enumerate}
                \item 意願
                \item 感情
                \item 判斷
                \item 在獨立句中使用表示祝願或第三人稱命令式: Qu'il parte tout de suite!
                \item 有些動詞在作肯定式使用時，從句動詞用陳述式，但否定或疑問式時，從句動詞用虛擬式
                    \begin{center}
                        Je crois qu'il \textbf{est} à Paris. Je ne crois pas qu'il \textbf{soit} à Paris.
                    \end{center}
            \end{enumerate}
        \subsection{Remark}
            當主從句的主語相同時，用不定式: Il veut que vous veniez vs. Il veut venir.
        
        \subsection{狀語從句}
            \begin{enumerate}
                \item 以下連詞引導的狀語從句需要虛擬式:
                    \begin{itemize}
                        \item avant que: \textbf{引導的從句中，動詞前往往加 ne}: Partez tout de suite avant qu'il ne soit trop tard.
                        \item bien que 雖然
                        \item quoique 儘管
                        \item sans que 無需
                        \item pour que 為了
                        \item afin que 以便
                        \item à condition que 只要
                        \item à moins que 除非
                        \item qui que 不論是誰
                        \item quoi que 不論是什麼
                        \item où que 不論哪裡
                        \item quel que 不論什麼樣的 \textbf{quel 要與修飾名詞性, 數一致} 
                    \end{itemize}
                \item 不用虛擬式:
                    \begin{itemize}
                        \item après que
                        \item parce que
                        \item depuis que 自從
                        \item pendant que 在$\cdots$期間
                    \end{itemize}
            \end{enumerate}
        
        \subsection{關係從句}
            先行詞所指的內容出於主觀願望或判斷, 從句要用虛擬式
    
    \section{虛擬式過去時 Le Subjonctif Passé}
        avoir être 虛擬式現在時 + 動詞過去分詞
        
        \subsection{用法}
            與現在時相同，但表示說話時已經完成或在將來某一時間前要完成的動作

    \section{條件式過去時 Le Conditionnel Passé}
        \subsection{構成}
            avoir (aur-), être (ser-) 的條件式現在時 + 動詞過去分詞
        \subsection{用法}
            \begin{enumerate}
                \item 設想，遺憾或惋惜等語氣
                \item 表示推測，即某事可能已發生，待進一步證實，常用於新聞
                \item 主從復合句中，表示過去未發生或無法實現的動作，僅為一種假設
                    \begin{itemize}
                        \item 從句: si + 愈過去時
                        \item 主句: 條件式過去時
                    \end{itemize}
                    \begin{center}
                        Si vous étiez venu hier, vous l'auriez recontré.
                    \end{center}
                    不要與現在時被動態混淆:
                    \begin{center}
                        S'il pleuvait, cette visite serait annulée (條件式現在時的被動態)
                        \\
                        S'il avait plu, nous aurions annulé cette visite (條件式過去時)
                    \end{center}
            \end{enumerate}

\chapter{造詞}
    \section{}{轉詞 La Conversion}
        \begin{enumerate}
            \item adj. $\to$ n. malade 患病的 $\to$ malade 病人
            \item 動詞不定式 v. $\to$ n. pouvoir 能夠 $\to$ le pouvoir 權利, vivre 生存 $\to$ les vivres 糧食
            \item 現在分詞 $\to$ 名词、形容词、介词: 
                \begin{itemize}
                    \item habitant $\to$ un habitant 居民
                    \item C'est une histoire \textbf{touchante} (形容詞)
                    \item Suivant cette théorie (介词)
                \end{itemize}
            \item 過去分詞 $\to$ 名詞、形容詞 reçu $\to$ un reçu, la fenêtre est \textbf{ouverte} toute la naquit
            \item 介詞 $\to$ 副詞: Il prit son fusil et s'en alla \textbf{avec} 他拿起步槍就走了
        \end{enumerate}

    \section{前綴 La Préfixation}
        \begin{enumerate}
            \item in-, im-, ir- 否定
            \item dé-, dés- 解除，分開
            \item re-, ré- 再次，重新
            \item sur- 超越，高於
            \item pré- 在$\cdots$之前
            \item co- 共同
            \item anti- 反抗，反對
            \item mal-, mé- 相反，不
            \item auto- 自身，自動
            \item télé- 遙遠的，電視的
        \end{enumerate}
            
    \section{後綴 Le Suffixation}
        \begin{enumerate}
            \item -er, -ir 動作 calmer, offrir
            \item -tion, -age 行為或動作結果
            \item -aire, -eur, ien 施動者，職業
            \item -té, -ité 品質
            \item -logie 學科
            \item -isme 學說，體系
            \item -able, -ible 可能性
            \item -ment 方式
        \end{enumerate}

    \section{復合構詞法 La Composition}
        \begin{enumerate}
            \item 以連詞符``-''相連: nouveau-né 新生兒
            \item 介詞 à, de 連結: fer à repasser 熨斗 pouvoir d'achat 購買力
            \item 無連結成分: prix plafond 最高價格 plein emploi 充分就業 chapeau melon 圓頂禮帽 gris perle 珠灰色
        \end{enumerate}

    \section{縮略構詞法 L'abreviation}
        \begin{enumerate}
            \item 縮減
            \item 首字母
            \item 合成 franglais (français + anglais)
        \end{enumerate}

\end{document}
